% Generated by simulations/python/examples.py -- do not edit by hand
\begin{table}[htbp]\centering\footnotesize
\setlength{\tabcolsep}{4pt}
\caption{Winners in the worked examples}\label{tab:examples}
\sbox{\tabbox}{\begin{tabular}{lrrrrrrrrr}
\toprule
Example & Plur. & Borda & Schul. & Fallb. & PVC & Unan. & Q 25\% & Q 33\% & Q 50\% \\
\midrule
E1 (compromise) & A & B & A & B & B & B & B & B & A \\
E2 (flooding) & A & A & A & A & A & B & B & B & A \\
E3 (cycle) & A & A & A & B & A & C (stuck) & B (fb) & A & A \\
E4 (status quo) & X & X & X & Y & X & S (stuck) & X & X & X \\
\bottomrule
\end{tabular}}
\usebox{\tabbox}
\tabnotes{E1: 51 voters rank A, B, C and veto C; 49 rank B, C, A and veto A. E2: E1 after the majority proposes two variants of A, ranked just below A by the majority and last by the minority, which vetoes them. E3: 35 voters rank A, B, C; 33 B, C, A; 32 C, A, B; each bloc vetoes its last option. E4: 99 voters rank X, Y, S and veto S; one voter ranks S, Y, X and vetoes X and Y. The status quo is C in E1--E3 and S in E4. Stuck: no proposal passes and the status quo stays; fb: every option is vetoed by the quota and QVC keeps the least-vetoed ones. Fallb.: fallback bargaining; PVC: Schulze within the proportional veto core; Unan.: unanimity with the status quo as default; Q $x$: QVC with a blocking quota $\theta = x$.}
{Profiles of Section~\ref{sec:examples}; \texttt{examples.py}.}
\end{table}
